% Interpretations for Chapter 4 (Unit 5). Numbers come from numbers.tex (generated from the notebook run);
% the confusion-matrix counts are read from Figure 4.8 (db_confusion.png).

\newcommand{\DBCVINTERP}{%
Every one of the ten ensembles beats both baselines on F1 and on ROC-AUC. Ensembles reach CV F1 0.788–0.806, against
0.730 for Logistic Regression and 0.723 for the Decision Tree; their ROC-AUC is 0.968–0.979, against 0.962 and 0.856.
\begin{itemize}
\item \textbf{Boosting} leads: Gradient Boosting (0.806), XGBoost (0.804), LightGBM (0.803) and CatBoost (0.802) all
reach ROC-AUC $\approx$0.979.
\item \textbf{Meta-ensembles} come next: Stacking (0.805) and Soft Voting (0.803) perform about as well as the best
single booster but do not beat it. Their four base models are all tree ensembles trained on the same features and make
highly correlated errors, so by Equation~(\ref{eq:var}) a large $\rho$ leaves little variance to remove.
\item \textbf{Bagging} ranks slightly lower on F1 (0.788–0.799) because it is very conservative: precision 0.91–0.99
with recall 0.66–0.69. AdaBoost is the extreme case, with precision 0.999 and recall 0.666.
\item \textbf{The Decision Tree} has the highest recall (0.742) but poor precision (0.705) and a ROC-AUC of only 0.856,
the classic sign of a single overfitted tree, which bagging was designed to fix.
\end{itemize}
The fold-to-fold standard deviation is $\le$0.007. The gaps among the top boosters ($\approx$0.003) are therefore within
noise, while the ensemble–baseline gap ($\approx$0.07–0.08) is roughly ten times larger.}

\newcommand{\DBTESTINTERP}{%
The held-out test results reproduce the CV ranking, so the models generalise. The best test F1 values are XGBoost
(0.813), Gradient Boosting (0.813), Soft Voting (0.811) and Random Forest (0.811), against 0.733 for Logistic Regression
and 0.727 for the Decision Tree.
\begin{itemize}
\item \textbf{Accuracy hides the gap.} Accuracy differs by only about one point (97.2\% vs 95.9\%) because 91\% of
patients are negative. F1 (+0.08) and PR-AUC (0.88 vs 0.81 and 0.56) reveal the real difference.
\item \textbf{Ranking quality.} Boosting reaches ROC-AUC 0.977–0.978 and PR-AUC 0.879–0.883.
\item \textbf{Training cost.} XGBoost and LightGBM train in about 1–2 s, against about 20 s for Gradient Boosting at the
same accuracy, while Stacking (about 50 s) adds cost without adding accuracy.
\item \textbf{Calibration.} AdaBoost ranks well (ROC-AUC 0.975), but its probabilities are poorly calibrated (Brier
0.110 vs 0.024 for the other boosters), because its output is a vote margin rather than a probability.
\end{itemize}}

\newcommand{\DBTHRINTERP}{%
The final model was chosen \emph{before} looking at the test set: \textbf{\dbBestName{}}, which has the highest CV F1.
Its out-of-fold F1-optimal threshold is \dbBestThr, almost exactly 0.5, giving test F1 \dbBestFone, precision
\dbBestPrec, recall \dbBestRec{} and ROC-AUC \dbBestAuc. Unlike the bank-marketing task in Chapter 2, threshold
optimisation brings \emph{no material gain} here: every model's test F1 changes by less than $\pm$0.008. Two facts
explain this. The boosted models are already well calibrated, and their F1 curve is flat over a wide range of thresholds
(Figure~\ref{fig:dbthr}), so 0.5 is already near-optimal. A lower threshold still trades precision for recall: for
example, the tuned LightGBM at 0.40 reaches recall 0.713 at precision 0.932. That trade is useful when screening, where
missing a diabetic costs more than a follow-up test. Hyperparameter tuning of LightGBM raised its CV PR-AUC to
\dbTuneAp{} but left its F1 essentially unchanged: the default boosting settings were already close to optimal for this
data.}

\newcommand{\DBROCINTERP}{%
In ROC space all ensembles lie close to the top-left corner. Logistic Regression is not far behind (0.960), because a
large part of the ranking is driven by HbA1c and glucose, which a linear model can also use. The precision–recall curves
separate the models much more clearly. Every ensemble keeps \textbf{precision 1.0 up to about 68\% recall}: these are
the patients with HbA1c $>$6.6\% or glucose $>$200, who are diabetic with certainty. Beyond that point the curves fall,
because the remaining diabetics have normal-range laboratory values. Logistic Regression's precision starts dropping
much earlier (around 40–45\% recall), and the single tree is a two-point line (PR-AUC 0.556). On this imbalanced problem,
PR-AUC is therefore the more informative ranking metric.}

\newcommand{\DBCMINTERP}{%
Of the 1,696 diabetic patients in the test set, Logistic Regression finds 1,083 and raises 176 false alarms.
\dbBestName{} finds \textbf{1,188 (105 more)} with only \textbf{40 false alarms (77\% fewer)}. At the tuned threshold
(0.49) the counts barely change (1,189 found, 45 false alarms). Out of 17,530 non-diabetics, only 40 are flagged, so
when the model says ``diabetic'' it is right 96.7\% of the time. The remaining 508 missed diabetics are the hard cases
identified in the EDA, with HbA1c and glucose in the normal range. They are the main limitation of any model trained on
these eight variables.}

\newcommand{\DBTHRCALINTERP}{%
\textbf{Left:} \dbBestName's test F1 stays at 0.80–0.81 for every threshold between about 0.35 and 0.95. Logistic
Regression peaks sharply near 0.5 (0.733) and falls quickly on either side. The ensemble is therefore both better and
more robust to the choice of threshold. \textbf{Right:} all three models are well calibrated, with points close to the
diagonal, so a predicted probability of 0.8 means roughly an 80\% observed rate. \dbBestName{} has the lowest Brier
score (0.0237, against 0.0249 for Random Forest and 0.0333 for Logistic Regression), so its probabilities can be
reported to clinicians directly as risk estimates.}

\newcommand{\DBSIZEINTERP}{%
These two curves show the ensemble mechanisms directly. \textbf{Boosting reduces bias:} after one stage the gradient
boosting model has test ROC-AUC \dbGbFirst; as stages are added, each fitting the residual errors of the ones before,
it rises to \dbGbLast. \textbf{Bagging reduces variance:} a single tree of the Random Forest scores ROC-AUC \dbRfOne;
averaging more trees raises this steadily to \dbRfAll{} with 300 trees. The curve flattens after about 100 trees, as
Equation~(\ref{eq:var}) predicts once the averaged variance approaches the correlated floor $\rho\sigma^2$.}

\newcommand{\DBSIGINTERP}{%
Using the same five folds for every model, the paired $t$-test shows that every ensemble's F1 gain over Logistic
Regression (+\dbSigMinGain{} to +\dbSigMaxGain) is statistically significant (largest $p$ = \dbSigMaxP{} $<$ 0.05).
The Decision Tree is the only model that is not better than Logistic Regression. On the test set, McNemar's test
comparing \dbBestName{} with Logistic Regression gives $\chi^2$ = \dbMcChi{} with $p = \dbMcP$:
\dbBestName{} alone classifies \dbMcBest{} patients correctly, against only \dbMcLr{} for Logistic Regression alone.
The improvement is not due to chance.}

\newcommand{\DBSHAPINTERP}{%
SHAP values \cite{shap} decompose each prediction into additive feature contributions, measured in log-odds, relative
to the average prediction. For \dbShapModel{}, the most influential features by mean $|\text{SHAP}|$ are
\dbShapTop. In the summary plot, high values (red) of HbA1c and blood glucose push risk strongly upwards, and older age,
higher BMI, hypertension and heart disease add further risk. The dependence plots show a sharp jump in contribution at
HbA1c $\approx$6.5\% and glucose $\approx$200 mg/dL. These are the diagnostic thresholds of the American Diabetes
Association \cite{ada}, which the model learned from data alone. Permutation importance on the test set agrees:
shuffling \dbPermFirst{} lowers PR-AUC by \dbPermFirstVal, \dbPermSecond{} by \dbPermSecondVal{} and \dbPermThird{} by
\dbPermThirdVal. Because the model's reasoning matches clinical knowledge, clinicians can audit and trust it more
readily.}

\newcommand{\DBCOMPARE}{%
\textbf{Same task, same test set, same metrics.} The direct comparison required by the review is between the ensembles
and the single-model baselines on identical folds and the same test set (Tables~\ref{tab:dbcv}–\ref{tab:dbtest}).
\dbBestName{} improves on Logistic Regression by +0.079 F1 (0.812 vs 0.733), +0.018 ROC-AUC and +0.067 PR-AUC, and on
the Decision Tree by +0.085 F1 and +0.115 ROC-AUC. All differences are statistically significant.

\textbf{Earlier units.} Logistic Regression and the Decision Tree are the same model families that were used as
baselines in Unit 2. The bank-marketing task in Chapter 2 uses a different dataset, so its numbers are not compared
directly. The same pattern appears, however: gradient-boosted trees beat linear models on tabular data with non-linear,
interacting effects. Choosing a threshold on out-of-fold predictions improved F1 there (0.487 $\rightarrow$ 0.531) but
not here, because the diabetes models are already well calibrated with a flat F1 curve. This shows that the threshold
step should be validated per task, not assumed to help.}

\newcommand{\DBDISCUSS}{%
\begin{itemize}
\item \textbf{Performance:} boosting ensembles (Gradient Boosting, XGBoost, LightGBM, CatBoost) are the most effective
classifiers for diabetes risk, with F1 $\approx$0.81 and ROC-AUC $\approx$0.98 on unseen patients. Bagging is close
but more conservative (high precision, lower recall). Meta-ensembles match but do not exceed the best booster.
\item \textbf{Generalisation:} CV and test results agree closely (CV F1 0.806 vs test 0.812 for the final model), and
fold standard deviations are $\le$0.007, so the results are stable and not an artefact of one split.
\item \textbf{Interpretability:} SHAP and permutation importance recover the clinical risk factors and thresholds. This
answers the ``black-box'' concern usually raised against ensembles.
\item \textbf{Training cost:} XGBoost and LightGBM train in 1–2 s on about 77,000 patients, much faster than
scikit-learn Gradient Boosting (about 20 s) or Stacking (about 50 s) at equal accuracy, which makes them the practical choice for
deployment.
\item \textbf{Limitations:} (i) about 30\% of diabetics have normal-range laboratory values and are still missed, so
the model should support, not replace, clinical judgement; (ii) HbA1c and glucose are recorded at a few discrete
levels, so performance on raw continuous clinical data may be lower; (iii) the dataset comes from a single public
source with limited provenance, and external validation on an independent cohort is needed before any clinical use;
(iv) ``No Info'' smoking history is informative in this dataset, but this pattern may not transfer to other hospitals.
\end{itemize}}

\newcommand{\DBCONCL}{%
Unit 5 implemented ten ensemble learners across the bagging, boosting and meta-ensemble families and evaluated them
against Logistic Regression and Decision Tree baselines, using 5-fold stratified CV, a held-out test set, F1, ROC-AUC
and PR-AUC, significance tests and SHAP. The final model, \dbBestName{} (selected by CV), identifies patients at risk of
diabetes with \textbf{F1 \dbBestFone, ROC-AUC \dbBestAuc{} and accuracy \dbBestAccPct}, with well-calibrated,
clinically interpretable predictions. This result answers the PBL III question: ensemble classifiers, boosting in
particular, are significantly more effective than single models at identifying individuals at risk of diabetes. Within
the integrated project, this unit confirms the Unit 2 finding that combining learners is the most reliable strategy
for tabular data.}
